\subsection{Critical Code and Cultural Software}
Digital audio tools run on code, and code carries cultural weight. Critical Code Studies (CCS) proposes that we “read code the way we read poetry,” not in a narrowly literary sense but by “situating the code as a cultural object of inquiry,” attending to its histories, uses, and interconnections with other code, hardware, and cultures \parencite[p.~31]{Marino_2020}. CCS emerged alongside software and platform studies as a way to supplement analyses that bracketed code. It insists that, instead of treating it as “meaningless math," we “approach it as a culturally situated sign system full of denotations and connotations,” exploring “the way its complex and unique sign systems make meaning as it circulates through varied contexts” \parencite[pp.19-33]{Marino_2020}. This commitment entails close reading and context: “every level of the code and hardware is a potential source of meaning,” to be interpreted with “the archaeologist’s attention to detail and meaning in context” \parencite[pp.32-48]{Marino_2020}.

As a throughline, CCS resonates with software studies’ argument that software should be treated as an object of inquiry by fields that have not historically “owned” it \parencite{Fuller_2008}. It also aligns with broader claims that code is social and historical. As Rita Raley puts it, “Code may in a general sense be opaque and legible only to specialists, much like a cave painting’s sign system, but it has been inscribed, programmed, written. It is conditioned and concretely historical” \parencite[as cited in][p.~41]{Marino_2020}. Adrian MacKenzie similarly writes, “Code is written and run within situated practices, vocabularies, and shared tools,” underscoring that any artifact “offers a glimpse of the cultures in which it was produced and circulated” \parencite[as cited in][pp.~41–42]{Marino_2020}. 

If we bring these insights into the sonic domain, by inspecting the practices of SID, we can hear that it centers the mappings between action and sound as sites where assumptions are concretized in the very ways systems are made to “listen” and respond. That is because these choices, choices about timbre, intelligibility, and noise, are just as much a design decision as they are technical specifications \parencite{Franinovic_Serafin_2013}. In practice, this means reading the code and paratexts of digital audio workstations, plugins, and audio recognition systems as cultural software whose defaults encode sonic norms about what counts as legible, desirable, or “clean.” As an example, a growing body of work on speech technologies illustrates these stakes. When development pipelines are organized around standardized norms, dialectal and cultural variation can be misread or marginalized; like those of African American English speakers. This provides an audible instance of culture traveling into code that “listens” and defines what is listened for \parencite{Cunningham_etAl_2025}.

\subsection{Procedural Rhetoric and Algorithmic Expression}
If code is cultural, then procedures are expressive. CCS maintains that “meaning grows out of the functioning of the code but is not limited to the literal processes the code enacts,” inviting critique of how operations of control and structures of data flows advance worldviews \parencite[p.~39]{Marino_2020}. Even small routines function as “the hub of a spoked wheel,” linking procedure to design principles, platform constraints, and social contexts \parencite[p.~48]{Marino_2020}. This method-esq bridge clarifies and even complements, to some degree, Ian Bogost’s notion of procedural rhetoric: to “make arguments with computational systems,” through the constraints and logics they enact, persuading by doing rather than by saying \parencite[p. 3]{Bogost_2007}. Extending (and at times complicating) this very same view, Steve Holmes emphasizes how procedural persuasion also operates through embodied practices, the “procedural habits” by which systems shape what users learn to do and perceive as normal \parencite[pp.64-67]{Holmes_2017}. If we read this all together, CCS and this dual account of procedural rhetoric/habits shows why the “exploration of code is never a neutral activity, free from an epistemology or a world view \parencite[p. 237]{Holmes_2017}." By this logic, the same \emph{must} hold true for the implementation of code as well: defaults for rhythm, tuning, or signal “clean-up” are procedural arguments about musical/sonic time, pitch, and voice, normalizing some sonic procedural habits (practices) while rendering others exceptional, abnormal, or perhaps even more labor-intensive in the end.

The arguments that have been laid out above are not confined to just the code in some form of isolation. CCS even cautions against treating code as self-sufficient; interpretation must attend to the contexts of execution, authorship, hardware, and paratexts \parencite[pp.~41, 46–48]{Marino_2020}. SID takes this further by foregrounding mappings between gesture and sound as both the aesthetic and cultural decisions of creators: how a system hears and responds is designed, not given \parencite{Franinovic_Serafin_2013}. In this light, algorithmic listening in audio software and speech systems is, in some form or another, an argument. Choices baked into drivers, DSP chains, and interface affordances co-produce what is easy to hear, do, and imagine. Raley’s reminder that code is “inscribed” and “concretely historical” and MacKenzie’s account of situated practices together underscore that procedural expression is a cultural practice before it is a technical one \parencite[as cited in][pp.~41–42]{Marino_2020}. 

As a brief sonic illustration, let's take a look at the disparities that appear when systems parse non-dominant dialects. In this case, the authors of this investigation were looking into the algorithmic pipeline for Automatic Speech Recognition (ASR) for African American English Speakers. As hinted at earlier in this section. This was in an effort to indicate how a pipeline’s procedures can tacitly argue that some voices are less legible due to there underlying procedures (and likely code). The point here is that it is not the metrics per se, but the way processes encode listening norms \parencite{Cunningham_etAl_2025}. Taken from Bogost’s vantage point, what a system makes effortless (and what it makes onerous) is the very substance of its rhetoric; Holmes adds that such persuasion often works through ingrained practices rather than explicit modeling alone \parencite{Bogost_2007,Holmes_2017}.

\subsection{Race After Technology and Discriminatory Design}
Given the above conceptual convergence. Those being, code as cultural implementation and procedural habits as arguments. We can begin to explain why the material consequences of design matter. CCS explicitly calls for “the need for interpreting these structures with an eye toward fundamental human concerns, concerning race, ethnicity, gender, and sexuality; ... [and] surveillance and control...,” precisely because programs underwrite systems “that do everything from facial recognition to tactical simulation” \parencite[p.~53]{Marino_2020}. Or, as Marino elsewhere puts it, “the walls of a computer do not remove code from the world but encode the world and human biases” \parencite[p.~33]{Marino_2020}. Ruha Benjamin names the social formation that results when these inscriptions formulate into infrastructures as the “New Jim Code,” whereby ostensibly neutral technologies reproduce and amplify racial inequalities through discriminatory design. IN response to this, her call for “abolitionist tools” attempts to move far past narrow fairness "fixes" toward reimagining the purposes, powers, and publics of technological systems themselves \parencite[154]{Benjamin_2019a}. If heard from  a more sonic perspective, algorithmic listening presents itself a crucial site of intervention: who and what systems are built to hear (and produce), and under what evaluative criteria? These are design questions with important political and cultural stakes \parencite{Franinovic_Serafin_2013}.

If we were to place these authors in conversation with one another, we can see that they chart a more complete conceptual arc. Fuller’s software studies widens the aperture so that software can be studied beyond engineering alone \parencite{Fuller_2008}; Marino’s CCS supplies a hermeneutic for the code level and demands that we track cultural imprints in procedure and context \parencite{Marino_2020}. Bogost clarifies that procedures persuade, which in the sonic sense, means defaults and pipelines argue for certain possible worlds of sound; while Holmes foregrounds the embodied, habitual dimensions of that persuasion \parencite{Bogost_2007,Holmes_2017}. Then in comes Benjamin, pushing the critique to structural outcomes, insisting that discriminatory design is neither accidental nor peripheral and that abolitionist design agendas are both necessary and possible \parencite{Benjamin_2019a}. SID translates these claims into the idiom of interactions with sound: by redesigning action–sound mappings and evaluation regimes, we can recompose how algorithms “listen,” and how their productions can begin shifting defaults and redistributing audibility \parencite{Franinovic_Serafin_2013}. In this, we end with brief empirical illustrations from speech recognition that beautifully, yet simply, underscores the point: when pipelines are built on narrow norms, audible harms follow \parencite{Cunningham_etAl_2025}. Given all this, what we find at stake are the cultural constitutions of the sonic and aural itself. What counts as signal, what is dismissed as noise, and whose sonic practices set the defaults.